% year: תשפ"ג
% type: מבחן
% semester: א
% moed: א
% number: 2א
% points: 10
% categories: taylor
% source: exams/מבחנים/תשפג/מועד א תשפג.pdf

\begin{question}
(10 נק') השתמשו בקירוב לינארי כדי להעריך את הביטוי $\sqrt[3]{28}$.
\end{question}

\begin{hint}
בחרו נקודה קרובה ל-$28$ שבה השורש השלישי ידוע.
\end{hint}

\begin{solution}
נגדיר $g(x)=\sqrt[3]{x}=x^{1/3}$ ונקודת בסיס $x_0=27$, שבה $g(27)=3$. הנגזרת:
\[ g'(x)=\frac{1}{3}x^{-2/3}=\frac{1}{3\sqrt[3]{x^2}},\qquad g'(27)=\frac{1}{3\cdot 9}=\frac{1}{27}. \]
הקירוב הלינארי (המשיק) בסביבת $x_0$: $g(x)\approx g(x_0)+g'(x_0)(x-x_0)$. עם $x=28$, $x-x_0=1$:
\[ \sqrt[3]{28}\approx 3+\frac{1}{27}\approx 3.037. \]
(הערך האמיתי הוא $3.0366\ldots$.)
\end{solution}
